\subsection*{略語}
\addcontentsline{toc}{subsection}{略語}

\par\smallskip\noindent\refstepcounter{table}\label{tab:ch04-02}\textbf{表 \thetable: 第04章 ソフトウェア構築 表02}\par\smallskip
\begin{longtable}{L{0.16\linewidth}L{0.32\linewidth}L{0.44\linewidth}}
\toprule
\textbf{略語} & \textbf{日本語} & \textbf{説明} \\
\midrule
API & アプリケーションプログラミングインタフェース & ライブラリ、サービス、フレームワークを使うための公開された入り口。 \\
CI & 継続的統合 & 変更を頻繁に統合し、自動ビルドと自動テストで確認する実践。 \\
COTS & 商用既製品 & 自作ではなく、市販または既存の製品として利用するソフトウェア部品。 \\
DSL & ドメイン固有言語 & 特定領域の問題を表しやすいように作られた言語。 \\
IDE & 統合開発環境 & 編集、ビルド、テスト、デバッグなどを統合した開発環境。 \\
MDA & モデル駆動アーキテクチャ & モデルから実装へ変換する考え方を中心にした開発方式。 \\
PIM & プラットフォーム非依存モデル & 特定の実装環境に依存しない解決策のモデル。 \\
PSM & プラットフォーム固有モデル & 実装環境の詳細を含むモデル。 \\
TDD & テスト駆動開発 & テストを先に書き、そのテストを通すようにコードを書く開発スタイル。 \\
WYSIWYG & 見たままが得られる方式 & 画面上の見た目が成果物の見た目に近い編集方式。 \\
\bottomrule
\end{longtable}
